\documentclass[a4paper,12pt]{article}

% ========== Кодировка и язык ==========
\usepackage[utf8]{inputenc}           % для pdflatex
\usepackage[T2A]{fontenc}             % кириллица
\usepackage[russian]{babel}           % русский язык

% Если используете XeLaTeX, закомментируйте строки выше и раскомментируйте:
% \usepackage{fontspec}
% \setmainfont{Times New Roman}       % или другой шрифт

% ========== Математические пакеты ==========
\usepackage{amsmath, amssymb, amsthm, mathtools}
\usepackage{enumitem}                 % гибкие списки
\usepackage{geometry}                 % настройка полей
\geometry{left=2cm, right=2cm, top=2cm, bottom=2cm}

% ========== Оформление заголовков ==========
\usepackage{titlesec}
\titleformat{\section}{\Large\bfseries\sffamily}{Задача \thesection}{1em}{}
\titleformat{\subsection}{\normalsize\bfseries}{\thesubsection)}{0.5em}{}

% ========== Пользовательские команды ==========
\newcommand{\task}[1]{\section*{#1}}          % для задач без автоматической нумерации
\newenvironment{solution}{\par\noindent\textbf{Решение.}\quad}{\hfill$\square$\par}
\newenvironment{remark}{\par\noindent\textbf{Замечание.}\quad}{}

% ========== Теоремы и определения ==========
\theoremstyle{definition}
\newtheorem{defn}{Определение}[section]
\theoremstyle{plain}
\newtheorem{thm}{Теорема}[section]
\newtheorem{lem}{Лемма}[section]
\newtheorem{prop}{Утверждение}[section]

% ========== Дополнительные обозначения ==========
\DeclareMathOperator{\Z}{\mathbb{Z}}
\DeclareMathOperator{\Q}{\mathbb{Q}}
\DeclareMathOperator{\R}{\mathbb{R}}
\DeclareMathOperator{\N}{\mathbb{N}}
\DeclareMathOperator{\Mobius}{\mu}    % функция Мёбиуса

\begin{document}

% ===== Заголовок =====
\begin{center}
    {\Huge\sffamily Решения к листочку по алгебре} \\[0.5cm]
    \large Ваше имя, группа, дата
\end{center}
\vspace{1cm}

% ===== Задача 1 =====
\section{}   % автоматическая нумерация
\textbf{а)} Докажите, что множества $\Z$ и $2\Z$ изоморфны как абелевы группы относительно сложения. Изоморфны ли они как кольца?

\textbf{б)} Изоморфны ли как абелевы группы $(\R, +)$ и $(\Q, +)$? $(\Q, +)$ и $(\Z, +)$?

\begin{solution}
    \textit{Здесь напишите решение.}
\end{solution}

% ===== Задача 2 =====
\section{}
\textbf{а)} Докажите, что любая подгруппа $(\Z, +)$ имеет вид $n\Z$, где $n \ge 0$.

\textbf{б)} Найдите обратный элемент к $[157]$ в кольце $\Z/235\Z$ и все целые решения уравнения $157x + 235y = 11$.

\begin{solution}
    % Ваше решение
\end{solution}

% ===== Задача 3 =====
\section{}
Убедитесь, что $\Z[\sqrt{5}] = \{ a + b\sqrt{5} \mid a,b\in\Z \}$ является коммутативным кольцом с единицей без делителей нуля (областью целостности). Покажите, что это кольцо не является факториальным.

\begin{solution}
    % Ваше решение
\end{solution}

% ===== Задача 4 =====
\section{}
\textbf{а)} Чему равно 17-е натуральное число, дающее остатки $2,4,6,8$ от деления на $5,9,11,14$?

\textbf{б)} Сколько решений в кольце $\Z/360\Z$ имеет уравнение $x^2 = 2$?

\textbf{в)} Сколько решений в кольце $\Z/n\Z$ имеет уравнение $x^2 = 1$?

\begin{solution}
    % Ваше решение
\end{solution}

% ===== Задача 5 =====
\section{}
Рассмотрим группу обратимых элементов $(\Z/n\Z)^\times$. Определим функцию Эйлера формулой $\varphi(n) = |(\Z/n\Z)^\times|$.

\textbf{а)} Докажите теорему Эйлера: $a^{\varphi(n)} = 1$ для любого $a \in (\Z/n\Z)^\times$.

\textbf{б)} Покажите, что $\varphi(n) = n(1 - p_1^{-1})\dots(1 - p_k^{-1})$, где $n = p_1^{\alpha_1}\dots p_k^{\alpha_k}$, $p_i$ --- различные простые числа.

\begin{solution}
    % Ваше решение
\end{solution}

% ===== Задача 6 =====
\section{}
Введём на множестве $X$ функций $f:\Z_{>0} \to \R$ две операции:
\begin{enumerate}[label=(\arabic*)]
    \item поточечное сложение: $(f+g)(n) = f(n)+g(n)$;
    \item свёртка Дирихле: $(f*g)(n) = \sum_{d|n} f(d)g\left(\frac{n}{d}\right)$.
\end{enumerate}

\textbf{а)} Докажите, что $(X,+,\*)$ --- коммутативное кольцо с единицей (какая функция является единицей этого кольца?). Опишите все обратимые элементы.

\textbf{б)} Найдите обратный элемент к функции
\[
\mathbf{1}: \Z_{>0} \to \R, \quad n \mapsto 1.
\]
Этот элемент называется функцией Мёбиуса.

\begin{solution}
    % Ваше решение
\end{solution}

% ===== Задача 7 =====
\section{}
\textbf{а)} Пусть $\xi$ --- первообразный корень по модулю простого $p>2$. Докажите, что $\xi$ или $\xi+p$ является первообразным корнем по модулю $p^2$.

\textbf{б)} Докажите, что если $\xi$ --- первообразный корень по модулю $p$ и $p^2$, то $\xi$ --- первообразный корень по модулю $p^k$ для любого $k>0$.

\textbf{в)} Докажите существование первообразного корня по модулю $2p^k$ для всех простых $p>2$ и всех $k\in\N$.

\textbf{г)} Покажите, что первообразный корень в $\Z/n\Z$ существует только для $n=2,4,p^k,2p^k$, где $p>2$ --- простое, $k\in\N$. Иными словами, группа $(\Z/n\Z)^\times$ циклическая только при таких $n$.

\textbf{д)} Найдите количество первообразных корней в $\Z/n\Z$. Если $a$ --- первообразный корень по модулю $n$, как выражаются все остальные первообразные корни?

\begin{solution}
    % Ваше решение
\end{solution}

\end{document}
